\section{Pretraining, and the fate of the founding hypothesis}
\label{sec:objective}

\paragraph{Objective.}
Standard JEPA plus one term:
\begin{equation}
\mathcal{L} = \operatorname{MSE}(\hat z, z_{\mathrm{tgt}})
+ \lambda\,\Omega(z_{\mathrm{ctx}})
\;[+\; \lambda_z \operatorname{SmoothL1}(\hat s, s)],
\end{equation}
with $\Omega$ = SIGReg \citep{balestriero2025lejepa} at $\lambda{=}1$ (16
random 1-D projections pushed toward Gaussian; VICReg control was
indistinguishable, 1.193 vs 1.183 average MASE, but SIGReg constrains the
whole distribution where a variance hinge accepts bimodal collapse). The
bracketed term is the ESJEPA arm ($\lambda_z{=}0.1$,
Section~\ref{sec:esjepa}). Two implementation lessons that cost real
results: compute variance statistics per patch position, not pooled
(pooled 0.9990 vs per-position 0.0000, positional diversity masking true
collapse), and never regularize a detached tensor (zero gradient, inflated
reported loss). AdamW \citep{loshchilov2019adamw}, peak LR
$3\times10^{-4}$, warmup 10\% then cosine, bf16, batch 512 per GPU on three
3090s. The LR ceiling and the short budgets are measurements, not habits: a
cosine calibrated for 40 epochs never decays in 3 (the run oscillated,
0.345--0.445, and its ``best checkpoint'' was a noise trough), and the full
corpus is exhausted in about one epoch, after which validation rises.
Collapse is monitored by per-position std and effective rank, calibrated
per corpus (rank at init: 3.7--4.8 on Monash, 43--87 on LOTSA).

\subsection{The founding bet, tested to destruction}
\label{sec:hypothesis}

The claim under test: latent extrapolation beats reconstruction as a
pretraining objective, increasingly so at long horizon. Four measurements,
in order.

\textbf{1. The paired duel is null.} One-variable reconstruction arm
(decode to normalized patches, SmoothL1, no anti-collapse: fixed targets
cannot collapse), checkpoints paired at 0.4\% val loss. Seven-dataset MASE:
latent 1.150, reconstruction 1.166. Reconstruction wins 17/28 paired cells,
median gap 0.5\% in its favor. GIFT-Eval: \gpair{0.979}{0.6766} vs
\gpair{1.003}{0.6764}, identical CRPS. Verdict: not ``JEPA loses'' but
``indistinguishable,'' which suffices to strike the claim.

\textbf{2. A horizon signal appeared and replicated.} Latent advantage
monotone in the GIFT term: $+0.8\%$ (short, $n{=}55$), $+3.1\%$ (medium,
$n{=}21$), $+6.6\%$ (long, $n{=}21$); same pattern in the Nixtla cells.
For a week this was the thesis.

\textbf{3. The per-step analysis killed the obvious mechanism.} If
predicting far is cheaper in latent space, the gap must grow with depth
\emph{inside} a window. Paired bootstrap over per-step MAE, eight datasets,
20{,}000 resamples: the gap \emph{shrinks} with depth, slope $-1.64\%$ per
100 steps, 95\% CI $[-2.82, -0.56]$, negative in 6/8 datasets. Across
rollout counts the picture inverts: gap vs number of rolls, Spearman
$+0.404$, $p{=}0.019$; latent significantly \emph{worse} at h192
(single window), ahead only at h720 (three rolls). Surviving candidate:
\emph{the latent objective degrades less under iteration on its own
outputs}. Fragile ($p{=}0.069$ without one dataset), with a registered
three-seed protocol not yet funded. A native h512 fine-tune (grown query
table, same pretrain, halving the roll count for the affected horizons)
supplied a paired check: on configurations unaffected by a window-induced
corpus change (Section~\ref{sec:results}), the multi-rollout group improved
by 1.7\% relative to the single-roll group, the sign the boundary
hypothesis predicts and the magnitude of a minor lever, consistent with the
flat term aggregates.

\textbf{4. What pretraining itself buys.} Against from-scratch at equal
fine-tuning: in-domain, nothing (1.193 vs 1.187, strict equality); on
never-fine-tuned datasets, $-26\%$ MASE on 8/8 (upper bound: the strongest
dataset was later found contaminated, Appendix~\ref{app:pitfalls}), tightened
to $-8.7\%$ in a budget-controlled clean rerun. Pretraining does not
improve fit; it improves transfer. And note the selector trap: fine-tuning
validation loss preferred the wrong model (0.159 scratch vs 0.181
pretrained) while the pretrained one transferred better.

One argument for JEPA survives untouched, and it is architectural, not
empirical: latents allow equivariance by objective. A context at
resolution $k_1$ and a target at $k_2$ can be tied through a predictor
conditioned on $w{=}k_2/k_1$; reconstruction cannot do this, two
resolutions do not share target values. Built (the FiLM path), never
funded. It is the cleanest reason to keep the formulation.
